% =============================================================================
% INCOHERENTIA: DE CONTINUATIONE (ARABIC ONLY)
% Status: DRAFTED
% Placement: Between XIX and XX
% Fully Arabic narrative chapter. Aisha, Zaynab, Saladin on the road.
% No translation provided. Readable to Arabic readers; structure available
% to others. Core themes: transmission, partial understanding, continuation.
% =============================================================================

\chapter*{{\AmiriFont الاستمرار}}
\addcontentsline{toc}{chapter}{{\AmiriFont الاستمرار}}

\setstretch{1.1}
{{\AmiriFont


ساروا في الصباح دون أن يعلن أحد بداية الرحلة.

لم يكن هناك أمر بالانطلاق، ولا إشارة واضحة تدل على الانتقال،
لكن الخطوات تتابعت، واحدة بعد الأخرى، حتى أصبح التقدّم
أمرًا مفروغًا منه.

كانت المدينة قد تراجعت خلفهم، لا كشيء ينتهي، بل كشيء
يفقد قدرته على التأثير.

لم تلتفت عائشة. لم يكن في ذلك امتناع، بل اختيار.

قالت زينب وهي تسير إلى جانبها:

«الناس هناك ما زالوا يتحدثون.»

أجابت عائشة:

«سيتحدثون دائمًا.»

«عن ماذا؟»

«عن ما يمكنهم تحمّله.»

سكتت زينب لحظة، ثم قالت:

«وهل يمكنهم تحمّل ما حدث؟»

لم تجب عائشة فورًا. نظرت إلى الطريق.

ثم قالت:

«يمكنهم تحمّل النسخة منه.»

كان صلاح الدين يسير خلفهما بقليل. قال:

«رأيت ما حدث.»

لم تلتفت عائشة، لكنها قالت:

«وأنت الآن تعرف.»

قال:

«المعرفة لم تغيّر النتيجة.»

قالت زينب:

«لكنها غيّرت ما لا يمكن إنكاره.»

سكتوا. لم يكن الصمت فراغًا، بل امتلاءً بما لا يحتاج إلى تكرار.

تابعوا السير.

الشمس ارتفعت ببطء، لا لتعلن يومًا جديدًا، بل لتكشف ما كان موجودًا بالفعل.

الطريق لم يتغيّر. لكن قدرتهم على قراءته تغيّرت.

وهذا كان كافيًا.

}}
